% Self-contained research fragment; no document preamble.
% All labels use gap:.  Sources: Differential Polygamma/Lerch, Polygamma
% Tail, and Bernoulli--Tornheim
% incoming reports at ProveIt 446dca9bd65c20e409d1d2561d4ffad408103e2a.

\section{The optimal moment prefix for every even-tail gap of four}
\label{gap:section}

The incoming \emph{Differential Polygamma and Lerch Identities} and
\emph{Polygamma Tail Identities} reports prove the optimal consecutive
moment prefix for adjacent even orders, $q-p=2$, and explicitly leave
the next gap $q-p=4$ open. The incoming
\emph{Bernoulli, Tornheim, and Digamma Identities} report proves all
pairs with $p=2,4,6$ and gives a residual determinant of size $p/2$
for arbitrary even pairs. We prove the entire gap-four family. The
first case outside those already proved families is $(p,q)=(8,12)$
\cite{DPL,PTI,BTD}.

The Differential report already reduces a fixed gap to a residual
determinant of size $(q-p)/2$, whose columns satisfy a positive
recurrence. We rederive that reduction with its normalization visible,
then prove the missing strict separation of column ratios. This shows
that the two-column determinant cannot vanish. These are statements
about rational coefficient matrices; no arithmetic independence of
zeta values is asserted. The fixed-gap input is
\cite[\texttt{ar:gap-det}, \texttt{ar:gap-factor}]{DPL}.

\subsection{The inherited moment coordinates and their conventions}

For even integers $p=2a$, $q=2b$, with $1\le a\le b$, write
\[
 A_p(n)=\zeta(p,n+1)=\zeta(p)-H_n^{(p)},\qquad x_n=n+\tfrac12,
\]
and let $D_r$ denote the continued value at $s=-2r$ of
\[
 D_{p,q}(s)=\sum_{n\ge0}A_p(n)A_q(n)x_n^{-s}.
\]
The source results prove regularity at these even arguments and the
inverse Bernoulli transforms
\begin{align}
 D_r&=\sum_{j=0}^r\frac{\binom{2r}{2j}}{2j+1}
 B_{2r-2j}(1/2)C_j,\label{gap:forward}\\
 C_j&=\sum_{r=0}^j\binom{2j+1}{2r}2^{-2(j-r)}D_r.
 \label{gap:inverse}
\end{align}
For $j\ge b$ their exact late coordinates are
\begin{equation}\label{gap:late-coordinates}
 C_j=\sum_{k=0}^{b-1}a_{jk}\zeta(2k),\qquad \zeta(0)=-\tfrac12.
\end{equation}
With $d=2j+1$ and $e=d-p-q+2k+1$, the inherited entries are zero
for $e<0$, and otherwise
\begin{equation}\label{gap:source-entry}
 a_{jk}=B_e\left\{
 \frac{\binom{d-q+1}{p-2k}}{d-q+1}
 +\frac{\binom{d-p+1}{q-2k}}{d-p+1}\right\}.
\end{equation}
Each binomial summand is omitted unless its lower index lies between
$2$ and its upper index. In particular, lower index zero is omitted.
The Bernoulli convention is $B_1=-1/2$.

Set
\begin{equation}\label{gap:late-matrix}
 M_{a,b}=(a_{b+r,k})_{0\le r,k<b}.
\end{equation}
In the formal coordinate order
\[
 C_0,\ldots,C_{b-1},\quad\zeta(0),\zeta(2),\ldots,\zeta(2b-2),
\]
the determinant of the complete first-$q$ moment matrix is
\begin{equation}\label{gap:full-prefix-determinant}
 \frac{\det M_{a,b}}{(4b-1)!!}.
\end{equation}
Indeed the first $b$ rows of the $C$ coordinate matrix form an
identity block, the last $b$ rows have the block $M_{a,b}$, and
the forward transform has diagonal $1/(2r+1)$. These facts and
\eqref{gap:forward}--\eqref{gap:source-entry} are source inputs;
the nonvanishing argument that follows is the new result.

\subsection{A complementary determinant of size half the gap}

Put $m=b-a\ge1$ and $h(t)=\coth(t/2)$. For an even polynomial
$Q(t)=\sum_{n=1}^b q_nt^{2n}$ define
\begin{equation}\label{gap:operator}
 \mathcal L_mQ(t)
 =\bigl(Q^{(2m)}(t)-Q^{(2m)}(0)\bigr)h(t)
                      +(Q(t)h(t))^{(2m)}.
\end{equation}
This is an odd formal Taylor series: the subtraction cancels the
possible pole in the first product, and $Qh$ is an odd Taylor series.
Let $K_{m,b}$ map $(q_1,\ldots,q_b)$ to the coefficients of
$t,t^3,\ldots,t^{2b-1}$ in \eqref{gap:operator}.

\begin{lemma}[Exact operator conjugacy]\label{gap:conjugacy}
If $M_{a,b}^{\mathrm{rev}}$ reverses the columns of $M_{a,b}$, then
\begin{equation}\label{gap:conjugacy-equation}
 K_{m,b}=
 \operatorname{diag}_{0\le r<b}\!\left(\frac2{(2r+1)!}\right)
 M_{a,b}^{\mathrm{rev}}
 \operatorname{diag}_{1\le n\le b}\!\bigl((2n)!\bigr).
\end{equation}
Consequently
\begin{equation}\label{gap:conjugacy-determinant}
 \det K_{m,b}=(-1)^{b(b-1)/2}4^bb!\det M_{a,b}.
\end{equation}
\end{lemma}

\begin{proof}
Write
\[
 h(t)=\sum_{k\ge0}h_kt^{2k-1},\qquad
 h_0=2,\quad h_k=\frac{2B_{2k}}{(2k)!}\quad(k\ge1).
\]
For use in all subsequent formulas define
\begin{equation}\label{gap:AB}
 A_r=\begin{cases}0,&1\le r\le m,\\
 (2r)^{\underline{2m}},&r\ge m+1,
 \end{cases}
 \qquad B_N=(2N-1)^{\underline{2m}}.
\end{equation}
The underline denotes a falling factorial. Row $r$ of $K_{m,b}$
corresponds to $N=r+m+1$. Its entry in column $n$ is
\begin{equation}\label{gap:operator-entry}
 h_{N-n}(A_n+B_N)\quad(n\le N),
\end{equation}
and is zero for $n>N$. The exceptional value $A_m=0$ comes from
the subtraction $Q^{(2m)}(0)$.

To compare this with the source, put $j=b+r$ and $k=b-n$ in
\eqref{gap:source-entry}. The Bernoulli index becomes $2(N-n)$,
while the two upper and lower binomial indices are
\[
 (2r+2,\,2n-2m),\qquad(2r+2m+2,\,2n).
\]
Multiplying by $2(2n)!/(2r+1)!$ gives exactly
\eqref{gap:operator-entry}. The first summand is omitted for
$n\le m$, precisely as required by \eqref{gap:AB}. This proves
\eqref{gap:conjugacy-equation}. The product of its positive diagonal
factors is
\[
 2^b\prod_{n=1}^b\frac{(2n)!}{(2n-1)!}=4^bb!,
\]
and column reversal has sign $(-1)^{b(b-1)/2}$.
\end{proof}

Put
\begin{equation}\label{gap:beta}
 \beta_k=\frac{2|B_{2k}|}{(2k)!}
        =\frac{4\zeta(2k)}{(2\pi)^{2k}}>0\quad(k\ge1),
\end{equation}
and, for $N\ge m+1$, put
\begin{equation}\label{gap:d}
 d_N=2\left\{(2N)^{\underline{2m}}+
                    (2N-1)^{\underline{2m}}\right\}>0.
\end{equation}
Define an array of rational numbers $C_{N,j}^{[m]}$, with
$N\ge1$ and $1\le j\le m$, by
\begin{align}
 C_{r,j}^{[m]}&=\delta_{rj}&& (1\le r\le m),
 \label{gap:initial}\\
 C_{N,j}^{[m]}&=\frac1{d_N}\sum_{r=1}^{N-1}
 \beta_{N-r}(A_r+B_N)C_{r,j}^{[m]}
 &&(N\ge m+1).
 \label{gap:recurrence}
\end{align}
Every entry $C_{N,j}^{[m]}$ with $N\ge m+1$ is strictly positive.
The superscript distinguishes this array from the inherited moment
coordinates $C_j$.

\begin{theorem}[Inherited fixed-gap boundary reduction]
\label{gap:complementary}
For $1\le a<b$ and $m=b-a$,
\begin{equation}\label{gap:complementary-formula}
 \det M_{a,b}
 =\frac{(-1)^{b(b-1)/2}}{4^bb!}
 \left(\prod_{N=m+1}^{b+m}d_N\right)
 \det\bigl(C_{b+i,j}^{[m]}\bigr)_{1\le i,j\le m}.
\end{equation}
Thus the remaining determinant has size $m=(q-p)/2$, independently
of the larger order $q$.
\end{theorem}

\begin{proof}
The sign relation $h_k=(-1)^{k-1}\beta_k$ for $k\ge1$ shows that
\begin{equation}\label{gap:formal-null-solutions}
 Q_j(t)=\sum_{n\ge1}(-1)^{n-j}C_{n,j}^{[m]}t^{2n}
 \qquad(1\le j\le m)
\end{equation}
formally satisfy $\mathcal L_mQ_j=0$. Conversely, the coefficient
at $t^{2N-2m-1}$ determines the next unknown coefficient uniquely:
its nonzero leading factor is $d_N$. Hence these are the unique
formal null solutions with their specified first $m$ coefficients.
This argument requires no convergence of the formal series.

Replace the first $m$ columns of $K_{m,b}$ by the corresponding
linear combinations specified by the truncations of
\eqref{gap:formal-null-solutions} through $n=b$. The column
transformation has determinant one, since those truncations have
initial block equal to the identity. Their first $b-m$ rows vanish.
Their bottom $m$ rows form the block $-TQ_{\mathrm{future}}$, where
$T$ is lower triangular with diagonal $d_{b+1},\ldots,d_{b+m}$,
and
\[
 (Q_{\mathrm{future}})_{ij}
      =(-1)^{b+i-j}C_{b+i,j}^{[m]}.
\]
This follows by writing the omitted terms in each exact formal
coefficient equation; row $b+i$ omits only coefficients
$b+1,\ldots,b+i$.

Move the first $m$ columns to the end. The other diagonal block is
triangular with diagonal $d_{m+1},\ldots,d_b$. Therefore
\[
 \det K_{m,b}
 =(-1)^{m(b-m)+m}\left(\prod_{N=m+1}^{b+m}d_N\right)
                          \det Q_{\mathrm{future}}.
\]
Factoring the row and column signs in $Q_{\mathrm{future}}$
gives $\det Q_{\mathrm{future}}=(-1)^{mb}
\det(C_{b+i,j}^{[m]})$. The combined exponent
$2mb-m(m-1)$ is even. Apply
\eqref{gap:conjugacy-determinant} to obtain
\eqref{gap:complementary-formula}.
\end{proof}

\subsection{Strict separation of the recurrence columns}

The classical Bernoulli Fourier formula, DLMF~24.8.1 at zero,
gives the second equality in \eqref{gap:beta} \cite{DLMF}. In particular, with
$x_\ell=(2\pi\ell)^{-2}$,
\[
 \beta_k=4\sum_{\ell\ge1}x_\ell^k\qquad(k\ge1).
\]
Strict Cauchy--Schwarz implies
\begin{equation}\label{gap:strict-log-convexity}
 \beta_{k+1}^2<\beta_k\beta_{k+2}\qquad(k\ge1),
\end{equation}
since the positive support contains distinct points. Thus
$\rho_k=\beta_{k+1}/\beta_k$ is strictly increasing. This familiar
moment argument is used solely to certify nonzero exact determinants.

\begin{lemma}[Ratio separation for every gap]\label{gap:ratio-lemma}
Let $m\ge2$. For any two initial columns $1\le i<j\le m$, the sequence
\begin{equation}\label{gap:ratio}
 R_N=\frac{C_{N,i}^{[m]}}{C_{N,j}^{[m]}}
       \qquad(N\ge m+1)
\end{equation}
is strictly increasing. Consequently every two-by-two minor of the
positive part of this array, with both its row indices and its
initial-column indices increasing, is strictly negative.
\end{lemma}

\begin{proof}
All relevant entries are positive. Since $A_1=\cdots=A_m=0$,
divide the recurrence by the second column's initial forcing
$B_N\beta_{N-j}/d_N$. Set
\begin{equation}\label{gap:normalized-weights}
 f_N=\frac{\beta_{N-i}}{\beta_{N-j}},\qquad
 w_{N,r}=\frac{\beta_{N-r}}{\beta_{N-j}}
                    \left(1+\frac{A_r}{B_N}\right).
\end{equation}
The exact ratio equation is
\begin{equation}\label{gap:ratio-equation}
 R_N=
 \frac{f_N+\displaystyle\sum_{r=m+1}^{N-1}
                   w_{N,r}C_{r,j}^{[m]}R_r}
      {1+\displaystyle\sum_{r=m+1}^{N-1}
                   w_{N,r}C_{r,j}^{[m]}}.
\end{equation}
By \eqref{gap:strict-log-convexity}, $f_N$ strictly increases: it is
a fixed-length product of consecutive $\rho_k$. For every fixed
$r\ge m+1>j$, the quotient $\beta_{N-r}/\beta_{N-j}$ is the
reciprocal of a fixed-length product of consecutive $\rho_k$,
so it strictly decreases as $N$ increases from $r+1$. Moreover
$A_r>0$ and $B_N$ strictly increases. Thus $w_{N,r}$ strictly
decreases with $N$ for each such fixed $r$.

At $N=m+1$ the sums are empty, giving $R_{m+1}=f_{m+1}$.
For the induction step assume the ratios through $R_{N-1}$ are
strictly increasing. Multiply \eqref{gap:ratio-equation} by its
denominator, subtract $R_{N-1}$ times that denominator, and subtract
the preceding row's corresponding zero identity. The term $r=N-1$
vanishes. The result is
\begin{align}
 &\left(1+\sum_{r=m+1}^{N-1}w_{N,r}C_{r,j}^{[m]}\right)
                 (R_N-R_{N-1})\notag\\
 &\quad=f_N-f_{N-1}
  +\sum_{r=m+1}^{N-2}(w_{N,r}-w_{N-1,r})C_{r,j}^{[m]}
                              (R_r-R_{N-1})>0.
 \label{gap:ratio-increment}
\end{align}
The first term is strictly positive. Each term in the sum is
nonnegative, because both differences there are negative. This
also handles the first induction step, when the sum is empty.
The assertion about minors follows by factoring out the positive
entries in their second column.
\end{proof}

\subsection{The gap-four determinant and the optimal prefix}

Specialize to $m=2$. Write
\[
 u_N=C_{N,1}^{[2]},\qquad v_N=C_{N,2}^{[2]}.
\]
For explicit calculation their common recurrence is
\begin{align}
 &(u_1,u_2)=(1,0),\qquad(v_1,v_2)=(0,1),\notag\\
 &y_N=\frac{\displaystyle\sum_{r=1}^{N-1}
                   \beta_{N-r}(A_r+B_N)y_r}
                {16(N-1)^2(2N-1)(2N-3)}\qquad(N\ge3),
 \label{gap:four-recurrence}\\
 &A_1=A_2=0,\qquad
 A_r=(2r)(2r-1)(2r-2)(2r-3)\quad(r\ge3),\notag\\
 &B_N=(2N-1)(2N-2)(2N-3)(2N-4).\notag
\end{align}
Define the strictly positive rational number
\begin{equation}\label{gap:delta}
 \Delta_b=u_{b+2}v_{b+1}-u_{b+1}v_{b+2}>0.
\end{equation}
Its sign follows immediately from Lemma~\ref{gap:ratio-lemma}.

\begin{theorem}[Every pair $q=p+4$]\label{gap:main}
For every integer $b\ge3$, let $p=2b-4$, $q=2b$. Then
\begin{equation}\label{gap:four-determinant}
 \det M_{b-2,b}
 =(-1)^{b(b-1)/2+1}
 \frac{4^b((b+1)!)^2(2b+3)!!(2b+1)!!}{3b!}\,\Delta_b\ne0.
\end{equation}
The moments $D_0,\ldots,D_{q-1}$ therefore have the full formal
rank $q$. Every later centered moment is an explicit rational
linear combination of these first $q$ moments. The determinant
of their full coordinate matrix is
\eqref{gap:four-determinant} divided by $(4b-1)!!$.
\end{theorem}

\begin{proof}
In \eqref{gap:complementary-formula} the endpoint determinant is
$-\Delta_b$. Here
\[
 d_N=16(N-1)^2(2N-1)(2N-3),
\]
so direct multiplication gives
\[
 \prod_{N=3}^{b+2}d_N
 =\frac{16^b((b+1)!)^2(2b+3)!!(2b+1)!!}{3}.
\]
This proves \eqref{gap:four-determinant}. Its scalar factor is
nonzero and $\Delta_b>0$, establishing nonvanishing. The full
rank assertion follows from \eqref{gap:full-prefix-determinant}.
To reconstruct a later moment, use \eqref{gap:inverse} to obtain
$C_0,\ldots,C_{2b-1}$ from the first $2b$ moments, invert
$M_{b-2,b}$ to recover all the even-zeta coordinates, and then use
\eqref{gap:late-coordinates} and \eqref{gap:forward}. Every step is
finite rational arithmetic.
\end{proof}

\subsection{A compact ordinary identity for the first new pair}

For $(p,q)=(8,12)$ the determinant and first inverse row are
\begin{align}\label{gap:812-data}
 \det M_{4,6}&=\frac{263589503}{309657600},\notag\\
 (M_{4,6}^{-1})_{0,*}
 &=\frac8{790768509}
 (469839,1269227,-11819495,56443363,83672540,12834702).
\end{align}
Thus the displayed integer row applied to $C_6,\ldots,C_{11}$
has value $-790768509/16$.

Define the odd polynomial
\begin{equation}\label{gap:812-K}
 \begin{split}
 K(N)=N^{13}\bigl(&469839+1269227N^2-11819495N^4
                   +56443363N^6\\
                 &+83672540N^8+12834702N^{10}\bigr)
 \end{split}
\end{equation}
and its centered difference
\begin{equation}\label{gap:812-W}
 W(x)=K(x+\tfrac12)-K(x-\tfrac12).
\end{equation}
The resulting ordinary, absolutely convergent identity is
\begin{equation}\label{gap:812-ordinary}
 \begin{split}
 \sum_{n=0}^{\infty}\biggl\{&W(x_n)A_8(n)A_{12}(n)
       -\frac{295198146}{77}x_n^4
       -\frac{10254604347}{154}x_n^2
       -\frac{59115956909}{385}\biggr\}
 =-\frac{790768509}{16}.
 \end{split}
\end{equation}
The summand is $O(x_n^{-2})$. To see this directly, the centered
Hurwitz expansion gives the first three terms of $A_8(n)A_{12}(n)$;
multiplication by the degree-$22$ polynomial $W$ gives precisely
the three displayed nondecaying terms. The inverse Bernoulli
transform identifies its continued sum with the integer combination
in \eqref{gap:812-data}; the removed even powers have continued
sum zero because $\zeta(-2r,1/2)=0$ for $r\ge0$. The accompanying
finite rational telescoping certificate in Section~\ref{tel:812}
supplies a second proof of this individual ordinary identity.

In polygamma notation the tail product is
\[
 A_8(n)A_{12}(n)
   =\frac{\psi^{(7)}(n+1)\psi^{(11)}(n+1)}{7!\,11!}.
\]
The compact difference definition \eqref{gap:812-W} avoids listing
twelve large coefficients while specifying the weight exactly.

\subsection{Verification, source update, and remaining limits}

The supplied exact script independently constructs the source matrix
from its binomial formula, constructs the operator matrix by symbolic
differentiation on monomials, and constructs the two recurrence
solutions. Its $323$ rational checks verify conjugacy for
$3\le b\le12$, the determinant formula for $3\le b\le18$,
formal null-series coefficients through $N=24$, and the finite sign
relations used in the proof. These computations are diagnostics for
finite formulas; the proof above establishes the infinite family.
A second standalone replay checks $45$ complementary determinants
for all $1\le a<b\le10$ and $140$ pairwise ratio minors at gaps
$m=2,3,4,5$. Thus the two tail scripts contain $508$ exact checks
in total; none extrapolates its finite range into a proof.

The optimal-prefix source conjecture can now record $q-p=4$ as
proved, alongside the inherited equal-order, adjacent-order, and
$p=2,4,6$ families. The fixed-gap reduction
\eqref{gap:complementary-formula} is the Differential report's
Proposition~\texttt{ar:gap-det}, rederived here; it retains that
attribution. The new input is Lemma~\ref{gap:ratio-lemma}, which
settles its two-dimensional boundary determinant. No new false
inherited theorem was established in this targeted audit.

For $m\ge3$, pairwise ratio separation does not by itself imply
nonvanishing of the $m$-square endpoint determinant. In fact its
sign need not remain fixed as $b$ changes. The next precise target
is the $m=3$ determinant, corresponding to $q-p=6$, with the
already proved overlapping pairs retained. Further useful targets
are a generating function for the two universal solutions in
\eqref{gap:four-recurrence}, and an exact recurrence or arithmetic
normalization for $\Delta_b$ itself. A general higher-minor sign
claim would require an additional proof; it is not a consequence
of the positivity used here.
